\begingroup

\newcommand{\autodiff}{algorithmic differentiation}
\newcommand{\Autodiff}{Algorithmic differentiation}
\newcommand{\TangentState}{u}
\newcommand{\ModelState}{\bm{u}}
% \newcommand{\CoreParameters}{\bm{y}}
% \newcommand{\HostParameters}{\hat{\bm{y}}}
% \newcommand{\AdjointState}{\eta}
\newcommand{\Objective}{J}
% \newcommand{\Residual}{\bm{g}}
\newcommand{\HostToCore}{Y}

\newcommand{\HostParameter}{\widehat{y}}
\newcommand{\HostParameters}{\widehat{\bm y}}
\newcommand{\CoreParameters}{\bm y}
\renewcommand{\State}{\bm u}
\newcommand{\HostSpace}{\widehat{\mathcal Y}}
\newcommand{\CoreSpace}{\mathcal Y}
\newcommand{\StateSpace}{\mathcal U}
\newcommand{\ResponseSpace}{\mathcal Z}
\newcommand{\ModelFamily}{\mathcal M}
\newcommand{\Residual}{\mathcal G}
\newcommand{\ParameterMap}{\Phi}
\newcommand{\StateMap}{S}
\newcommand{\Observation}{Q}
\newcommand{\ResponseMap}{R}
\newcommand{\ComposedResponse}{\widehat{R}}
\renewcommand{\Objective}{J}
\newcommand{\AdjointState}{\bm\eta}
\newcommand{\D}{\mathrm D}
\newcommand{\Tan}[1]{\operatorname{Tan}\left[\,#1\,\right]}

\newcommand{\haty}{\hat{\bm{y}}}
\newcommand{\corey}{\bm{y}}
% \newcommand{\ModelState}{\bm{u}}
\newcommand{\residual}{\bm{g}}
\newcommand{\nhost}{\hat{n}}
\newcommand{\ncore}{n}
\newcommand{\nstate}{m}

% Maps
\newcommand{\Ymap}{\ParameterMap} %\mathcal{Y}}   % host -> core parameter map
\newcommand{\Mmap}{\mathcal{M}}   % core parameters -> state (FE solve)
\newcommand{\Tmap}{\ComposedResponse_{\ModelFamily}}%\mathcal{T}}   % composed tangent map


\section{Section Formulation}

After the primal section state has been converged (Box~\ref{box:impl}), one differentiates the enhanced equilibrium equation at fixed generalized deformation \(\bm{e}\) to recover the enhanced sensitivity \(\bm{\eta}_{,\HostParameter}\).
The strain and stress sensitivities at each integration point then follow from the differentiated kinematic map and the material sensitivity relation. 
The section sensitivity is obtained by integrating the differentiated generalized stress-resultant operator.

\begin{SaveEquation}{eq:section-sensitivity-condensation}
\Knn \bm{\eta}_{,\HostParameter} = -\residn_{,\HostParameter}^{\mathrm{dir}}
\end{SaveEquation}
where \(\residn_{,\HostParameter}^{\mathrm{dir}}\) is the derivative of the enhanced residual \eqref{eq:def-enhan-residual} with respect to the active host parameter \(\HostParameter\) while holding \(\bm{\eta}\) fixed. 



Differentiating the enhanced residual \eqref{eq:def-enhan-residual} furnishes:
\begin{SaveEquation}{eq:enhanced-residual-sensitivity-full}
\residn_{,\HostParameter}
=
\int_{\Section}
\left(
{\EnhanShapes}_{,\HostParameter}^{\top}\StressVector
+
\EnhanShapes^{\top}\StressVector_{,\HostParameter}
+
a_h\,\EnhanShapes^{\top}\StressVector
\right)\dA.
\end{SaveEquation}
where 
\[
a_h \coloneq \frac{(\dA)_{,\HostParameter}}{\dA}
\]
is the relative sensitivity of the section measure and the stress sensitivity \(\StressVector_{,\HostParameter}\) at each integration point \(\bm{r}_{\ell}\) is
\begin{SaveEquation}{eq:section-stress-sensitivity}
\StressVector_{,\HostParameter}
=
\StressVector_{,\HostParameter}^{\mathrm{mat}}
+ \mathbb{C}\StrainVector_{,\HostParameter}.
\end{SaveEquation}
with
\begin{SaveEquation}{eq:section-strain-sensitivity}
\StrainVector_{,\HostParameter}
=
\hat{\StrainVector}_{,\HostParameter}
+ {\EnhanShapes}_{,\HostParameter}\bm{\eta}
+ \EnhanShapes\,\bm{\eta}_{,\HostParameter}
\end{SaveEquation}
%
Using \eqref{eq:section-strain-sensitivity} and \eqref{eq:section-stress-sensitivity}
% \[
% \StressVector_{,\HostParameter}
% =
% \StressVector_{,\HostParameter}^{\mathrm{mat}}
% +
% \mathbb{C}\StrainVector_{,\HostParameter},
% \qquad
% \StrainVector_{,\HostParameter}
% =
% {\TotalStrainMatrix}_{,\HostParameter}\bm{e}
% +
% {\EnhanShapes}_{,\HostParameter}\bm{\eta}
% +
% \EnhanShapes\,\bm{\eta}_{,\HostParameter},
% \]
one obtains
\[
\residn_{,\HostParameter}
=
\Knn \bm{\eta}_{,\HostParameter}
+
\residn_{,\HostParameter}^{\mathrm{dir}}, 
\]
with \(\Knn\) furnished by \eqref{eq:Knn} and 
\begin{SaveEquation}{eq:enhanced-residual-sensitivity-direct}
\residn_{,\HostParameter}^{\mathrm{dir}}
=
\int_{\Section}
\Big[
{\EnhanShapes}_{,\HostParameter}^{\top}\StressVector
+
\EnhanShapes^{\top}
\bigl(
\StressVector_{,\HostParameter}^{\mathrm{mat}}
+
\mathbb{C}\bigl({\TotalStrainMatrix}_{,\HostParameter}\bm{e}
+ {\EnhanShapes}_{,\HostParameter}\bm{\eta}\bigr)
\bigr)
+
a_h\,\EnhanShapes^{\top}\StressVector
\Big]\dA.
\end{SaveEquation}

Hence the condensed sensitivity equation is
\[
\Knn \bm{\eta}_{,\HostParameter} = -\residn_{,\HostParameter}^{\mathrm{dir}}.
\]

% For the generalized resultants,
% \begin{SaveEquation}{eq:section-resultant-sensitivity-split}
% \bm{s}_{,\HostParameter}
% =
% \int_{\Section}
% \TotalStrainMatrix^{\top}\StressVector_{,\HostParameter}\,\dA
% +
% \bm{s}_{,\HostParameter}^{\mathrm{geom}},
% \end{SaveEquation}
% where the explicit operator and measure contribution is
% \begin{SaveEquation}{eq:section-resultant-geom}
% \bm{s}_{,\HostParameter}^{\mathrm{geom}}
% =
% \int_{\Section}
% \left(
% {\TotalStrainMatrix}_{,\HostParameter}^{\top}\StressVector
% +
% a_h\,\TotalStrainMatrix^{\top}\StressVector
% \right)\dA.
% \end{SaveEquation}

The full resultant derivative is then:
\[
\bm{s}_{,\HostParameter}
=
\int_{\Section}
\left(
{\TotalStrainMatrix}_{,\HostParameter}^{\top}\StressVector
+
\TotalStrainMatrix^{\top}\StressVector_{,\HostParameter}
+
a_h\,\TotalStrainMatrix^{\top}\StressVector
\right)\dA.
\]

\begin{ambox}[Section Sensitivity Determination.]{box:section-sensitivity}
\paragraph{Objective.}
Given generalized deformations \(\bm{e}\), the converged section state \(\bm{\eta}\), and an active parameter \(\HostParameter\), compute the conditional sensitivities \(\bm{\eta}_{,\HostParameter}\) and \(\bm{s}_{,\HostParameter}\) over the cross-section discretization.

\paragraph{Procedure.}
\begin{enumerate}
    \item Retrieve the converged section state \(\bm{\eta}\), together with the stresses \(\StressVector_\ell\) and material tangents \(\mathbb{C}_\ell\) at the quadrature points of the section discretization.

    \item Solve for \(\bm{\eta}_{,\HostParameter}\):
    \[
    \Knn \bm{\eta}_{,\HostParameter} = -\residn_{,\HostParameter}^{\mathrm{dir}}
    \]
    where
    \RepeatEquation{eq:enhanced-residual-sensitivity-direct}
    \item Integrate the section resultant sensitivity:
    \[
    \bm{s}_{,\HostParameter}
    =
    \int_{\Section}
    \TotalStrainMatrix^{\top}\StressVector_{,\HostParameter}\,\dA
    +
    \bm{s}_{,\HostParameter}^{\mathrm{geom}},
    \]
    \begin{itemize}

    \item At each quadrature point, evaluate the strain sensitivity from the differentiated kinematic relation:
    \[
    \StrainVector_{,\HostParameter}
    =
    \hat{\StrainVector}_{,\HostParameter}
    +
    {\EnhanShapes}_{,\HostParameter}\bm{\eta}
    +
    \EnhanShapes\,\bm{\eta}_{,\HostParameter}.
    \]

    \item Evaluate the stress sensitivity with the material sensitivity relation:
    \[
    \StressVector_{,\HostParameter}
    =
    \StressVector_{,\HostParameter}^{\mathrm{mat}}
    +
    \mathbb{C}\StrainVector_{,\HostParameter}.
    \]
    \end{itemize}

\end{enumerate}
\end{ambox}


\section{Element}

To accommodate parameters that may affect the element length, reference geometry, or interpolation data, it is convenient to express the element resisting force on a fixed parent domain \(\hat e\):
\begin{equation}
\bm{p}_a(\ModelState,\HostParameter)
=
\int_{\hat e}
J(\hat{\reflenx},\HostParameter)\,\bm{B}_a(\ModelState,\HostParameter)\,\bm{\sigma}(\ModelState,\HostParameter)\,\dd\hat{\reflenx}
=
\int_{\hat e}
J\,\bm{B}_a\,\bm{A}_*\,\bm{s}\,\dd\hat{\reflenx} .
\label{eq:pa-parent}
\end{equation}
Here \(\ModelState\) denotes the vector of discrete configuration variables in the chosen chart, \(J\) is the parent-to-element Jacobian, and \(\bm{\sigma}=\bm{A}_*\bm{s}\).

Define
\begin{equation}
(\,\cdot\,)_{,\HostParameter} \coloneq \frac{\dd}{\dd \HostParameter}(\,\cdot\,),
\qquad
(\,\cdot\,)_{;\HostParameter}
\coloneq
\left.\frac{\partial}{\partial \HostParameter}(\,\cdot\,)\right|_{\ModelState},
\label{eq:def-direct-h}
\end{equation}
so that \((\,\cdot\,)_{;\HostParameter}\) denotes the explicit parameter derivative at fixed discrete state. Differentiating \eqref{eq:pa-parent} along the equilibrium path gives
\begin{equation}
\bm{p}_{a,\HostParameter}
=
\int_{\hat e}
\left[
D_\chi[\bm{B}_a\bm{A}_*\mid\bm{s}] \cdot \bm{u}_h
+
J\,\bm{B}_a\bm{A}_*\,\bm{s}_{,\HostParameter}
+
\bigl(J\,\bm{B}_a\bm{A}_*\mid\bm{s}\bigr)_{;\HostParameter}
\right]\,\dd\hat{\reflenx} ,
\label{eq:pa-h-raw}
\end{equation}
where \(\bm{u}_h\) is the configuration sensitivity field induced by the nodal sensitivities,
\begin{equation}
\bm{u}_h^h
=
\sum_b \bm{\varphi}_{u b}\,\ModelState_{b,h}.
\label{eq:uh-field}
\end{equation}

The first term in \eqref{eq:pa-h-raw} is the state-induced geometric contribution. 
By definition of the discrete geometric operator,
\begin{equation}
D_\chi[\bm{B}_a\bm{A}_*\mid\bm{s}] \cdot \bm{u}_h
=
\sum_b \bm{G}_{ab}[\bm{s}]\,\ModelState_{b,h}.
\label{eq:geom-state-sens}
\end{equation}
The constitutive term is treated by differentiating the section response. The strain sensitivity admits the decomposition
\begin{equation}
\bm{e}_{,\HostParameter}
=
D\bm{e}\cdot \bm{u}_h + \bm{e}_{;\HostParameter}
=
\sum_b \bm{B}_b^\star \ModelState_{b,h} + \bm{e}_{;\HostParameter},
\label{eq:eh-split}
\end{equation}
and therefore
\begin{equation}
\bm{s}_{,\HostParameter}
=
\bm{C}_s\,\bm{e}_{,\HostParameter} + \bm{s}_{;\HostParameter}
=
\bm{C}_s\sum_b \bm{B}_b^\star \ModelState_{b,h}
+
\bm{C}_s\,\bm{e}_{;\HostParameter}
+
\bm{s}_{;\HostParameter}.
\label{eq:sh-split}
\end{equation}
Substituting \eqref{eq:geom-state-sens} and \eqref{eq:sh-split} into \eqref{eq:pa-h-raw} yields
\begin{equation}
\bm{p}_{a,\HostParameter}
=
\sum_b
\left[
\int_{\hat e}
J\Bigl(
\bm{B}_a\,\bm{C}_\sigma\,\bm{B}_b^\star
+
\bm{G}_{ab}[\bm{s}]
\Bigr)\,\dd\hat{\reflenx}
\right]\ModelState_{b,h}
+
\bm{p}_{a;\HostParameter},
\label{eq:pa-h-final}
\end{equation}
with
\begin{equation}
\bm{p}_{a;\HostParameter}
=
\int_{\hat e}
\left[
\bigl(J\,\bm{B}_a\bm{A}_*\mid\bm{s}\bigr)_{;\HostParameter}
+
J\,\bm{B}_a\bm{A}_*
\Bigl(
\bm{C}_s\,\bm{e}_{;\HostParameter}
+
\bm{s}_{;\HostParameter}
\Bigr)
\right]\,\dd\hat{\reflenx} .
\label{eq:pa-h-direct}
\end{equation}
The first term in \eqref{eq:pa-h-direct} expands to
\begin{equation}
\bigl(J\,\bm{B}_a\bm{A}_*\mid\bm{s}\bigr)_{;\HostParameter}
=
J_{,\HostParameter}\,\bm{B}_a\bm{\sigma}
+
J\,(\bm{B}_a\bm{A}_*)_{;\HostParameter}\,\bm{s}.
\label{eq:pa-h-direct-expand}
\end{equation}

Equation \eqref{eq:pa-h-final} gives the component sensitivity of the resisting force vector. In matrix form,
\begin{equation}
\bm{p}_{,\HostParameter}
=
\bm{k}\,\ModelState_{,\HostParameter}
+
\bm{p}_{;\HostParameter},
\label{eq:p-h-matrix}
\end{equation}
with element tangent
\begin{equation}
\bm{k}_{ab}
=
\int_{\hat e}
J\Bigl(
\bm{B}_a\,\bm{C}_\sigma\,\bm{B}_b^\star
+
\bm{G}_{ab}[\bm{s}]
\Bigr)\,\dd\hat{\reflenx} .
\label{eq:k-ab-sens}
\end{equation}
Thus the force sensitivity splits into a state-induced term, obtained by applying the consistent tangent to the nodal sensitivity vector, and a direct term that collects all explicit parameter dependence of the element geometry, interpolation, and constitutive response.

% For a pure material parameter, such as \(E\) or \(G\), the geometry and interpolation are fixed, so that
% \[
% J_{,\HostParameter}=0,
% \qquad
% (\bm{B}_a\bm{A}_*)_{;\HostParameter}=0,
% \qquad
% \bm{e}_{;\HostParameter}=0,
% \]
% and \eqref{eq:pa-h-direct} reduces to
% \begin{equation}
% \bm{p}_{a;\HostParameter}
% =
% \int_{\hat e}
% J\,\bm{B}_a\bm{A}_*\,\bm{s}_{;\HostParameter}\,\dd\hat{\reflenx}.
% \label{eq:pa-h-material}
% \end{equation}
% If the constitutive law is linear in the form \(\bm{s}=\bm{C}_s(h)\bm{e}\), then
% \begin{equation}
% \bm{s}_{;\HostParameter}=\bm{C}_{s,h}\bm{e},
% \qquad
% \bm{p}_{a;\HostParameter}
% =
% \int_{\hat e}
% J\,\bm{B}_a\bm{A}_*\,\bm{C}_{s,h}\bm{e}\,\dd\hat{\reflenx}.
% \label{eq:pa-h-linear-material}
% \end{equation}

At the assembled level, if the equilibrium residual is written as
\begin{equation}
\bm{g}(\ModelState,\HostParameter)=\bm{p}(\ModelState,\HostParameter)-\bm{f}(\ModelState,\HostParameter),
\label{eq:residual-sens}
\end{equation}
then differentiation gives
\begin{equation}
\bm{g}_{,\HostParameter}
=
\bm{K}\,\ModelState_{,\HostParameter}
+
\bm{p}_{;\HostParameter}
-
\bm{f}_{,\HostParameter},
\label{eq:residual-sens-final}
\end{equation}
where \(\bm{K}\) is the assembled consistent tangent formed from \eqref{eq:k-ab-sens}. 
This is the form required in direct differentiation procedures for nonlinear static and dynamic sensitivity analysis.


\endgroup
